\textbf{Curriculum length and ordering control.} At the main cell we varied the curriculum length $T_c$ and replaced the operand-size order by a random order with the same pacing (Table~\ref{tab:tc}, Fig.~\ref{fig:sweep}a; 5 seeds each; $T_c=1000$ is the main-group row).

\begin{table}[t]\centering\caption{Pacing sweep and shuffled-order control at the main cell (5 seeds). \emph{final acc} is held-out accuracy when the run stopped (at the 95\% crossing or at 4000 steps).}\label{tab:tc}
\footnotesize\begin{tabular}{lccc}
\toprule
Variant & steps\_to\_95 & reached & final acc \\
\midrule
curriculum, $T_c=250$ & 2905 $\pm$ 727 & 5/5 & 0.96 \\
curriculum, $T_c=500$ & 3320 $\pm$ 893 & 3/5 & 0.84 \\
curriculum, $T_c=1000$ & 3955 $\pm$ 101 & 1/5 & 0.75 \\
curriculum, $T_c=2000$ & 3990 $\pm$ 22 & 1/5 & 0.59 \\
curriculum, $T_c=4000$ & 4000 $\pm$ 0 & 0/5 & 0.54 \\
shuffled order, $T_c=1000$ & 4000 $\pm$ 0 & 0/5 & 0.49 \\
uniform (main group) & 2945 $\pm$ 1020 & 3/5 & 0.86 \\
\bottomrule
\end{tabular}
\end{table}

\begin{figure*}[t]\centering\includegraphics[width=0.75\textwidth]{sweep_grid.pdf}
\caption{(a) Steps to 95\% against curriculum length $T_c$ (black dots: seeds; lines: means, uniform and shuffled-order references). (b) Fraction of seeds reaching 95\% in each weight-decay/training-fraction cell (3 seeds; 5 at 1/0.5).}\label{fig:sweep}\end{figure*}

Shorter curricula are better: with $T_c=250$ the curriculum reaches 95\% in 5 of 5 seeds in 2905 steps on average, close to uniform (2945, 3 of 5), and the means degrade monotonically as $T_c$ grows, down to 0 of 5 reached at $T_c=4000$. The two longest curricula are confounded with the fixed 4000-step budget: the pool is not full until step $T_c$ and the training set is fitted only at about 1875 ($T_c=2000$) and 3700 ($T_c=4000$) steps, so little time is left for generalisation and part of the 0/5 and 1/5 outcomes is by construction. The clean contrast is $T_c=250$ against $T_c=1000$. The shuffled-order run with $T_c=1000$ reached 95\% in 0 of 5 seeds against 1 of 5 for the operand-size order. With 5 seeds this one-seed difference is not evidence for or against a benefit of ordering by operand size; what the sweep does show is that, within the curriculum family, a slower pool growth reached 95\% less often (5 of 5 at $T_c=250$ against 1 of 5 at $T_c=1000$). Relative to uniform sampling, however, no pacing or ordering was better, and the 7-seed check above does not establish that any of them is worse. A plausible reason (we did not isolate it) is that a small early pool is memorised quickly and the model then needs the remaining examples before it can find a solution that works for all pairs, so delaying full exposure delays generalisation; the training-accuracy column (about 960 steps to fit the training set under both orderings vs.\ about 125 for uniform) is consistent with this but does not prove it. We did not test an intermediate-sized operand-size curriculum with a different difficulty measure.

\textbf{Weight decay and training fraction.} The grid (Table~\ref{tab:grid}) is the second sensitivity study. Reading it against H3 and H4: generalisation within 4000 steps needed both a nonzero weight decay and a large enough training fraction, and in this region the ordering made at most a small difference compared with these two knobs. At $f=0.7$ the anti-curriculum was slower than uniform at $\lambda=1$ (2575 vs.\ 1817) and at $\lambda=3$ (2358 vs.\ 1142, with 2 of 3 seeds reaching 95\%), which is the clearest ordering effect in the grid, but with 3 seeds it remains suggestive.
